\chapter{Euclidean Sobolev inequalities and their endpoints}\label{sec:sobolev-foundations}
The local estimates for elliptic equations use an integrability gain that comes from Euclidean geometry. We prove that input here, starting with iterated Holder inequalities. The source is \cite{SimonPDE}, Lecture 17, Lemma 2 and Corollary 1, printed pp.~200--202. Throughout, $n\geq2$ and $W^{1,p}_0(\Omega)$ denotes the closure of $C_c^\infty(\Omega)$ in the $W^{1,p}$ norm. No regularity of the open set $\Omega\subset\mathbb R^n$ is required by that definition or by the zero-boundary inequalities below. Basic completion, weak derivatives, and Holder and Tonelli inequalities are background tools.

\begin{lemma}[An iterated Holder inequality for coordinate projections]\label{lem:coordinate-projection-holder}
For $i=1,\ldots,n$, let $F_i$ be a nonnegative integrable function of the $n-1$ coordinates $\widehat x_i$ obtained by omitting $x_i$. Then
\[
 \int_{\mathbb R^n}\prod_{i=1}^n F_i(\widehat x_i)^{1/(n-1)}\,dx
 \leq\prod_{i=1}^n\left(\int_{\mathbb R^{n-1}}F_i\right)^{1/(n-1)}.
\]
\end{lemma}
\begin{proof}
For $n=2$ this is the product identity from Tonelli. Suppose it is known in dimension $n-1$. Integrate in $x_n$ first and apply Holder to the first $n-1$ factors, with all exponents equal to $n-1$. Put $G_i(\widehat x_i,\widehat x_n)=\int_{\mathbb R}F_i(\widehat x_i)\,dx_n$ for $i<n$, where the argument of $G_i$ consists of the remaining $n-2$ coordinates. The integral in the statement is at most
\[
 \int_{\mathbb R^{n-1}}F_n(x')^{1/(n-1)}
             \prod_{i<n}G_i(\widehat x_i,\widehat x_n)^{1/(n-1)}\,dx'.
\]
Holder with exponents $n-1$ and $(n-1)/(n-2)$ bounds this by
\[
 \|F_n\|_1^{1/(n-1)}
 \left(\int_{\mathbb R^{n-1}}\prod_{i<n}G_i^{1/(n-2)}\,dx'\right)^{(n-2)/(n-1)}.
\]
Apply the induction hypothesis to the integral and use $\|G_i\|_1=\|F_i\|_1$. This gives the result. Zero factors are immediate; truncation to bounded functions on bounded boxes and monotone convergence justify the argument without any initial finiteness assumption on the product integral.
\end{proof}

\begin{theorem}[The first-power Sobolev inequality]\label{thm:sobolev-first-power}
For every $u\in C_c^1(\mathbb R^n)$,
\[
 \|u\|_{L^{n/(n-1)}}
 \leq\prod_{i=1}^n\|\partial_i u\|_{L^1}^{1/n}
 \leq\|Du\|_{L^1}.
\]
This is the $p=1$ input in \cite{SimonPDE}, Lecture 17, Lemma 2, pp.~200--201; no sharp constant is claimed.
\end{theorem}
\begin{proof}
On every coordinate line, the fundamental theorem of calculus and compact support give
\[
 |u(x)|\leq F_i(\widehat x_i),\qquad
 F_i(\widehat x_i)=\int_{\mathbb R}|\partial_i u(x)|\,dx_i.
\]
Multiply these $n$ inequalities and raise the result to the power $1/(n-1)$. Lemma \ref{lem:coordinate-projection-holder} gives
\[
 \int |u|^{n/(n-1)}\leq\prod_i\|\partial_i u\|_1^{1/(n-1)}.
\]
Taking the power $(n-1)/n$ proves the first inequality. Each $\|\partial_i u\|_1$ is at most $\|Du\|_1$, proving the second. If one coordinate derivative vanishes identically, the same argument forces $u=0$, consistently with compact support.
\end{proof}

\begin{theorem}[Sobolev gain below the critical derivative exponent]\label{thm:euclidean-sobolev-gain}
Let $1\leq p<n$ and $p^*=np/(n-p)$. For every open $\Omega\subset\mathbb R^n$ and every $u\in W^{1,p}_0(\Omega)$,
\[
 \|u\|_{L^{p^*}(\Omega)}\leq C_{n,p}\|Du\|_{L^p(\Omega)}.
\]
One may take $C_{n,1}=1$ and $C_{n,p}=(n-1)p/(n-p)$ for $p>1$. These constants are convenient bounds, not optimal values. See \cite{SimonPDE}, Lecture 17, Lemma 2.
\end{theorem}
\begin{proof}
First take $u\in C_c^\infty(\Omega)$ and extend it by zero. The case $p=1$ is Theorem \ref{thm:sobolev-first-power}. For $p>1$, put $\alpha=(n-1)p/(n-p)>1$ and apply that theorem to $|u|^\alpha\in C_c^1(\mathbb R^n)$. The chain rule and Holder give
\[
 \|u\|_{p^*}^{\alpha}
 \leq\alpha\int |u|^{\alpha-1}|Du|
 \leq\alpha\|u\|_{p^*}^{\alpha-1}\|Du\|_p,
\]
since $\alpha n/(n-1)=(\alpha-1)p/(p-1)=p^*$. If $u$ is nonzero, divide by the finite power of its norm; otherwise the conclusion is immediate.

For general $u\in W^{1,p}_0(\Omega)$ choose compactly supported smooth $u_j\to u$ in $W^{1,p}$. Their zero extensions converge in $W^{1,p}(\mathbb R^n)$, since both function and gradient norms of differences are unchanged by extension at this stage. The inequality for $u_j-u_k$ makes them Cauchy in $L^{p^*}$. The $L^{p^*}$ limit agrees almost everywhere with the zero extension of $u$, by taking subsequences of the $L^p$ and $L^{p^*}$ convergences. Pass to the limit. This argument uses the closure definition and does not assume a trace theorem for a rough boundary.
\end{proof}

\begin{corollary}[A volume-dependent zero-boundary Poincare estimate]\label{cor:volume-poincare}
If $0<\mathcal L^n(\Omega)<\infty$ and $1\leq p<n$, then
\[
 \|u\|_p\leq C_{n,p}\mathcal L^n(\Omega)^{1/n}\|Du\|_p
 \qquad(u\in W^{1,p}_0(\Omega)).
\]
\end{corollary}
\begin{proof}
Holder gives $\|u\|_p\leq\mathcal L^n(\Omega)^{1/p-1/p^*}\|u\|_{p^*}$. The exponent difference is $1/n$; apply Theorem \ref{thm:euclidean-sobolev-gain}. The zero-boundary space matters: the constant function one belongs to $W^{1,p}(\Omega)$ and has zero gradient, so the estimate cannot hold on all of that space when $\Omega$ has positive finite volume.
\end{proof}

\begin{theorem}[Every finite integrability gain in dimension two]\label{thm:two-dimensional-sobolev}
Let $\Omega\subset\mathbb R^2$ be open with $0<\mathcal L^2(\Omega)<\infty$. For every $2\leq q<\infty$ there is $C_q<\infty$ such that
\[
 \|u\|_{L^q(\Omega)}\leq C_q\mathcal L^2(\Omega)^{1/q}\|Du\|_{L^2(\Omega)}
 \qquad(u\in H_0^1(\Omega)).
\]
The constants may depend on $q$. This includes the finite-power assertion in \cite{SimonPDE}, Lecture 17, Corollary 1, p.~202.
\end{theorem}
\begin{proof}
Set $p=2q/(q+2)$, so $1\leq p<2$ and $p^*=q$. Approximation and Holder on the finite-volume domain imply $H_0^1(\Omega)\subset W^{1,p}_0(\Omega)$. Theorem \ref{thm:euclidean-sobolev-gain} and Holder for the gradient yield
\[
 \|u\|_q\leq C_{2,p}\|Du\|_p
 \leq C_{2,p}\mathcal L^2(\Omega)^{1/p-1/2}\|Du\|_2.
\]
The exponent difference is $1/q$. This proves the result, including $q=2$ through the first-power inequality. Letting $q\to\infty$ in the displayed assertion would require uniform control of its constants and is not justified.
\end{proof}

\begin{corollary}[The scaled gain on balls]\label{cor:scaled-sobolev-ball}
Let $\chi=n/(n-2)$ for $n>2$ and $\chi=2$ for $n=2$. For $R>0$ and $w\in H_0^1(B_R(x_0))$,
\[
 \left(\int_{B_R(x_0)}|w|^{2\chi}\right)^{1/\chi}
 \leq C_n R^{2-n+n/\chi}\int_{B_R(x_0)}|Dw|^2.
\]
The radius factor is one when $n>2$ and is $R$ when $n=2$.
\end{corollary}
\begin{proof}
On $B_1$, use Theorem \ref{thm:euclidean-sobolev-gain} with $p=2$ for $n>2$, and Theorem \ref{thm:two-dimensional-sobolev} with $q=4$ for $n=2$, then square the inequality. For $v(y)=w(x_0+Ry)$, the left side scales by $R^{n/\chi}$ and the gradient integral by $R^{n-2}$. Rearranging gives precisely the displayed factor. Smooth approximation extends the change of variables to the stated spaces.
\end{proof}

\begin{example}[Concentration forbids a larger Sobolev exponent]\label{ex:sobolev-exponent-obstruction}
For $1\leq p<n$ and finite $q>np/(n-p)$, no constant $C$ can make $\|u\|_q\leq C\|Du\|_p$ hold for every $u\in C_c^\infty(B_1)$. The same failure holds with the full $W^{1,p}$ norm on the right.
\end{example}
\begin{proof}
Choose a nonzero $\phi\in C_c^\infty(B_1)$ and set $u_\varepsilon(x)=\varepsilon^{1-n/p}\phi(x/\varepsilon)$, $0<\varepsilon<1$. Direct change of variables gives
\[
 \|Du_\varepsilon\|_p=\|D\phi\|_p,\qquad
 \|u_\varepsilon\|_p=\varepsilon\|\phi\|_p,\qquad
 \|u_\varepsilon\|_q=\varepsilon^{1-n/p+n/q}\|\phi\|_q.
\]
The last exponent is negative exactly when $q>np/(n-p)$, so that norm diverges while the right-hand norms remain bounded. The critical exponent is thus forced by concentration, independently of the particular proof of the inequality.
\end{proof}

\begin{example}[An unbounded function at the critical derivative exponent]\label{ex:critical-sobolev-unbounded}
For every $n\geq2$ there is an unbounded function in $W^{1,n}_0(B_1\subset\mathbb R^n)$. In particular, $H_0^1(B_1\subset\mathbb R^2)$ does not embed in $L^\infty$.
\end{example}
\begin{proof}
Choose a smooth radial cutoff $\eta$, equal to one on $B_{1/4}$ and zero outside $B_{1/2}$, with $0\leq\eta\leq1$. For $x\ne0$ put
\[
 u(x)=\eta(x)\log\log(e/|x|),
\]
and assign any finite value at zero. This function is unbounded in essential supremum near zero. Its $L^n$ norm is finite by polar integration. Near zero its gradient has length $1/(r\log(e/r))$, and
\[
 \int_0^{1/4}\frac{dr}{r\log(e/r)^n}<\infty
\]
because $n>1$. Integration by parts off $B_\delta$ shows these classical derivatives are weak derivatives: the inner boundary term is bounded by a constant times $\delta^{n-1}\log\log(e/\delta)$, which tends to zero.

For large $k$, cap the radial logarithmic function at height $k$ before multiplying by $\eta$. The resulting compactly supported functions are Lipschitz, agree with $u$ off a small ball, and converge to $u$ in $W^{1,n}$ by dominated convergence of their values and gradients. Their piecewise gradients are weak gradients because the capped functions are continuous across the matching sphere. Each can be approximated in $W^{1,n}$ by smooth compactly supported mollifications inside $B_1$, using translation continuity in $L^n$ for the function and its weak derivatives. Hence $u\in W^{1,n}_0(B_1)$. The endpoint failure does not contradict any finite-$q$ assertion in Theorem \ref{thm:two-dimensional-sobolev}.
\end{proof}
